% =============================================================================
% Chapter 4: Methodology
% =============================================================================

\chapter{Methodology}
\label{ch:methodology}

This chapter describes the experimental pipeline used to investigate whether behavioural feedback improves LLM-driven algorithm design. The pipeline comprises three stages, LLM screening, behavioural feature analysis and selection, and feedback format screening. We first describe the LLaMEA framework and benchmark configuration that underpin all three stages.

% ---------------------------------------------------------------------------
\section{LLaMEA Framework}
\label{sec:meth-llamea}

All experiments use the Large Language Model Evolutionary Algorithm (LLaMEA)~\cite{vanStein2024LLaMEA}, an evolutionary framework that treats an LLM as a generative engine for optimisation algorithms. Vanilla LLaMEA operates a $(1 \stackrel{+}{,} 1)$-ES loop. In each generation, the LLM produces a candidate algorithm as a Python class, the algorithm is evaluated on a benchmark suite, and the resulting fitness score is fed back to the LLM together with the algorithm's source code. The process is summarised in \cref{alg:llamea}.

\begin{algorithm}[t]
\caption{$(1+1)$-LLaMEA evolutionary loop, adapted from \cite{vanStein2024LLaMEA}.}
\label{alg:llamea}
\SetKwInOut{Input}{Input}
\SetKwInOut{Output}{Output}
\Input{Task prompt $S$, feedback prompt template $F_0$, budget $T$}
\Output{Best algorithm $a_b$ and its quality $y_b$}
$a_0 \gets \text{LLM}(S)$ \tcp*{Generate initial algorithm}
$(y_0, \sigma_0, e_0) \gets f(a_0)$ \tcp*{Evaluate}
$a_b \gets a_0;\; y_b \gets y_0$\;
\For{$t \gets 1$ \KwTo $T$}{
  $F \gets (S,\; \{(\text{name}(a_i), y_i)\}_{i=0}^{t},\; (a_b, y_b, \sigma_b),\; F_0)$\;
  $a_{t} \gets \text{LLM}(F)$ \tcp*{Mutate: refine or explore}
  $(y_{t}, \sigma_{t}, e_{t}) \gets f(a_{t})$\;
  \If{$e_{t} \neq \varnothing$}{$y_{t} \gets 0$ \tcp*{Runtime error: worst fitness}}
  \If{$y_{t} \geq y_b$}{$a_b \gets a_{t};\; y_b \gets y_{t};\; \sigma_b \gets \sigma_{t}$}
}
\Return{$a_b, y_b$}
\end{algorithm}

\begin{figure*}[!t]
  \centering
  \includegraphics[width=\textwidth]{figures/fig_mabbob_instances.pdf}
  \caption{Per-function weight distribution of the selected 10 MA-BBOB instances, coloured by function group. All 24 BBOB functions are covered with near-uniform balance.}
  \label{fig:mabbob-instances}
\end{figure*}

\begin{table*}[!t]
\centering
\caption{BBOB function weights for the 10 selected MA-BBOB instances. Only non-zero contributions are shown (as percentage of total weight). Functions are grouped: separable (f1--f5), low/moderate conditioning (f6--f9), high conditioning/unimodal (f10--f14), multimodal adequate (f15--f19), multimodal weak (f20--f24).}
\label{tab:instance-weights}
\footnotesize
\begin{tabular}{@{}r|ccccc|cccc|ccccc|ccccc|ccccc@{}}
\toprule
& \multicolumn{5}{c|}{\textbf{Sep.}} & \multicolumn{4}{c|}{\textbf{Low/mod}} & \multicolumn{5}{c|}{\textbf{High/uni}} & \multicolumn{5}{c|}{\textbf{MM adeq.}} & \multicolumn{5}{c}{\textbf{MM weak}} \\
\textbf{Inst.} & \footnotesize 1 & \footnotesize 2 & \footnotesize 3 & \footnotesize 4 & \footnotesize 5
& \footnotesize 6 & \footnotesize 7 & \footnotesize 8 & \footnotesize 9
& \footnotesize 10 & \footnotesize 11 & \footnotesize 12 & \footnotesize 13 & \footnotesize 14
& \footnotesize 15 & \footnotesize 16 & \footnotesize 17 & \footnotesize 18 & \footnotesize 19
& \footnotesize 20 & \footnotesize 21 & \footnotesize 22 & \footnotesize 23 & \footnotesize 24 \\
\midrule
191 & & & & & & 28 & & & & & & & & & & & & & 25 & 21 & 1 & & 25 & \\
277 & 15 & & 14 & & & & 10 & & & 13 & & & 13 & 6 & & & & 20 & & & & 6 & & 4 \\
300 & & & 19 & & & & 20 & & & & & 18 & & & & & & 12 & & & 29 & & 2 & \\
412 & & 10 & & & 12 & & & 13 & & 14 & & & & & & & 11 & & 11 & & & & & 29 \\
455 & & 13 & & & & & & 17 & & 14 & 14 & & & & & & & 14 & & & & & & 28 \\
635 & & & & 23 & 4 & & 13 & & & & & 9 & & 16 & & 23 & & & & & & 14 & & \\
648 & 5 & & & 9 & & & 20 & & 18 & & & & 18 & & 17 & & & & & 11 & 2 & & & \\
744 & & 17 & & & & 19 & & & & 10 & & & & 7 & 17 & 15 & 15 & & & & & & & \\
760 & & 1 & & 9 & 12 & & & 4 & 20 & & 7 & & 15 & & & & & & & & 17 & 14 & & \\
843 & & 21 & & 14 & & & & 14 & & 10 & 2 & 20 & & & & & 18 & & & & & & & \\
\bottomrule
\end{tabular}
\\[0.3em]
{\scriptsize Values are percentages of total instance weight. Empty cells indicate zero weight.}
\end{table*}

\paragraph{Initialisation.}
Every evolutionary run starts from a shared initial algorithm, \texttt{RandomSearch}, that performs uniform random sampling across the search space and returns the best sample found. Pre-seeding the population ensures that all models and conditions begin from the same starting point, so observed differences are attributable to the LLM and feedback rather than initialisation luck. LLaMEA's \texttt{initialize()} method skips LLM generation when the population is pre-filled~\cite{vanStein2024LLaMEA}.

\paragraph{Mutation prompts.}
Following the configuration identified as optimal by van Stein et al.~\cite{vanStein2025Behaviour}, we use a 90/10 refine/explore mutation mix. With 90\% probability the LLM receives a \emph{refine} prompt (``\textit{Refine and simplify the selected algorithm to improve it}''), and with 10\% probability an \emph{explore} prompt (``\textit{Generate a new algorithm that is different from those tried before}''). This balances exploitation of promising algorithm structures with exploration of novel designs.

\paragraph{Selection.}
We use elitist $(1+1)$ selection, where the offspring replaces the parent only if it achieves equal or better fitness. Recent work has explored larger population sizes. In a behaviour space study by van Stein et al.~\cite{vanStein2025Behaviour} compare six LLaMEA configurations including population-based variants with $\mu/\lambda = 4/12$, and the SAGE experiments~\cite{vanStein2026LLaMEASAGE} use $4$ parents and $16$ offspring. Notably, the $(1+1)$ elitist variant (LLaMEA-4) outperformed the population-based alternatives in the behaviour space study~\cite{vanStein2025Behaviour}. We retain the $(1+1)$ strategy for this reason and because the scale of our experiments (10 models $\times$ 5 seeds in Stage~1, 29 conditions $\times$ 5 seeds in Stage~3) makes larger populations computationally prohibitive.

\paragraph{Feedback.}
After each evaluation, the LLM receives a feedback string reporting the algorithm's performance. The standard (vanilla) feedback reports only the AOCC score and its standard deviation across evaluation runs. Our experiments augment this string with behavioural feature information.

% ---------------------------------------------------------------------------
\section{Benchmark Configuration}
\label{sec:meth-benchmark}

\paragraph{Problem suite.}
Candidate algorithms are evaluated on MA-BBOB (Many-Affine BBOB)~\cite{vermetten2023mabbob}, which constructs diverse problem instances by combining affine transformations of classical BBOB functions~\cite{hansen2021coco}. Each MA-BBOB instance is defined by a weight vector over the 24 BBOB functions, producing landscapes that blend multiple function characteristics. All instances are 5-dimensional with search bounds $[-5, 5]^5$.

\paragraph{Instance selection.}
From the 1{,}000 available MA-BBOB instances, we select 10 training instances using a greedy optimisation procedure that maximises function coverage and group balance (\cref{fig:mabbob-instances}). The selection criteria are: (1) all 24 BBOB functions must be represented, (2) each of the five function groups should contribute approximately 20\% of the total weight, and (3) per-function weight should be as uniform as possible. The resulting selection covers all 24 functions with group shares within $\pm$0.6\% of the ideal 20\%. \Cref{tab:instance-weights} shows the per-instance weight breakdown; each instance blends between 5 and 9 base functions.

\paragraph{Evaluation protocol.}
Each candidate algorithm is evaluated on all 10 instances across 5 random seeds, yielding 50 independent runs per candidate. The evaluation budget per run is $B = 2000 \times d = 10{,}000$ function evaluations, following the BLADE default~\cite{vanStein2025BLADE}. All evaluation is managed by IOHexperimenter~\cite{deNobel2024iohexperimenter}, which enforces budget limits, logs full trajectories, and ensures reproducibility.

\paragraph{Performance metric.}
Algorithm quality is measured by the Area Over the Convergence Curve (AOCC), an anytime performance metric in $[0, 1]$~\cite{vanStein2025BLADE}. For a single run of algorithm $a$ on function $f$ with budget $B$, the per-run area is:
\begin{equation}
\label{eq:aoc-run}
\text{AOC}(y_{a,f}) = \frac{1}{B} \sum_{i=1}^{B} \left(1 - \frac{\min(\max(y_i, lb), ub) - lb}{ub - lb}\right)
\end{equation}
where $y_i$ is the log-scaled current-best fitness at evaluation $i$, and $ub = 10^{2}$, $lb = 10^{-8}$ are the clipping bounds. The AOCC for a candidate algorithm is the average over all $N$ runs (instances $\times$ seeds):
\begin{equation}
\label{eq:aocc}
\text{AOCC}(a) = \frac{1}{N} \sum_{f,s} \text{AOC}(y_{a,f,s})
\end{equation}
In our setup $N = 50$ (10 instances $\times$ 5 seeds). This single scalar is reported to the LLM as feedback.

\paragraph{Behavioural feature computation.}
For every successful candidate evaluation, the full trajectory (all $T = 10{,}000$ evaluations with decision variables and fitness values) is logged via IOHexperimenter. From each trajectory, all 32 behavioural features described in \cref{ch:features} are computed post-hoc. Features are averaged across the 50 runs per candidate to produce a single feature vector per algorithm. This dataset, which comprises feature vectors for every valid candidate across all models and seeds, forms the basis for the Stage~2 feature analysis.

% ---------------------------------------------------------------------------
\section{LLM Screening}
\label{sec:meth-screening}

Before investigating feedback mechanisms, a suitable LLM must be identified. Small language models vary substantially in their ability to generate valid, high-quality optimisation code, and inference cost matters when each evolutionary run requires 100 LLM calls.

\paragraph{Models.}
Ten models are evaluated, spanning locally served open-weight models and cloud API models. \cref{tab:models} lists all models with their parameter counts and providers.

\begin{table}[t]
\centering
\caption{LLMs evaluated in Stage~1. Local models are served via Ollama~\cite{ollama}, while API models are accessed through Google's Gemini API. Parameter counts reflect total (not active) parameters.}
\label{tab:models}
\footnotesize
\begin{tabular}{@{}llrl@{}}
\toprule
\textbf{Model} & \textbf{Provider} & \textbf{Params} & \textbf{Backend} \\
\midrule
Qwen3.5 4B   & Qwen/Alibaba~\cite{qwen3}          & 4B  & Ollama \\
Qwen3.5 9B   & Qwen/Alibaba~\cite{qwen3}          & 9B  & Ollama \\
Qwen3.5 27B  & Qwen/Alibaba~\cite{qwen3}          & 27B & Ollama \\
RNJ-1 8B     & Essential AI~\cite{rnj1}            & 8B  & Ollama \\
Devstral Small~2 & Mistral AI~\cite{devstral2}     & 24B & Ollama \\
OLMo~3 7B    & Allen AI~\cite{olmo3}               & 7B  & Ollama \\
OLMo~3 32B Think & Allen AI~\cite{olmo3}           & 32B & Ollama \\
Granite~4 Micro  & IBM~\cite{granite4}             & 3B  & Ollama \\
\midrule
Gemini~3 Pro     & Google~\cite{gemini3flash}      & --- & API \\
Gemini~3 Flash   & Google~\cite{gemini3flash}      & --- & API \\
\bottomrule
\end{tabular}
\end{table}

\paragraph{Protocol.}
Each model runs the LLaMEA loop for 100 candidates across 5 independent seeds (seeds 0--4), all using vanilla AOCC-only feedback. The shared initial algorithm (RandomSearch) ensures identical starting conditions. All 32 behavioural features are computed for every successful candidate, producing a dataset of several thousand algorithm--feature observations.

\paragraph{Ranking criteria.}
Models are ranked on a composite of three criteria:
\begin{enumerate}
  \item \textbf{Failure rate}: the fraction of candidates that fail to produce valid fitness (due to syntax errors, interface mismatches, or runtime exceptions).
  \item \textbf{Best AOCC}: the highest AOCC achieved across all 5 seeds, measuring peak algorithm quality.
  \item \textbf{Inference efficiency}: tokens generated per second, which determines wall-clock time per evolutionary run.
\end{enumerate}
The top-performing model on this composite ranking is selected for Stage~3.

% ---------------------------------------------------------------------------
\section{Feature Analysis and Selection}
\label{sec:meth-feature-selection}

The Stage~1 dataset, which comprises behavioural feature vectors for all valid candidates across 10 models and 5 seeds, is analysed to determine which of the 32 features are most predictive of algorithm quality and to select a tractable subset for the feedback experiments.

\paragraph{Ranking methods.}
Three complementary approaches are applied:
\begin{enumerate}
  \item \textbf{Spearman rank correlation}~\cite{spearman1904proof} between each feature and AOCC, computed across all valid candidates. This captures monotonic associations regardless of their functional form.
  \item \textbf{Random Forest permutation importance}~\cite{breiman2001randomforests}: a Random Forest regression model is trained to predict AOCC from all 32 features, and feature importance is measured by the decrease in out-of-bag $R^2$ when each feature is permuted.
  \item \textbf{Kolmogorov--Smirnov distributional effect size}: the KS statistic between feature distributions of top-25\% and bottom-25\% algorithms by AOCC, measuring how well each feature separates high- from low-quality algorithms.
\end{enumerate}

\paragraph{Consensus ranking.}
Each method produces a ranked list of 32 features. These are aggregated via Borda count: each feature receives $32 - r_i$ points from each method (where $r_i$ is its rank in method $i$), and total scores determine the consensus ranking.

\paragraph{Redundancy pruning.}
Features are processed in Borda-rank order. If a candidate feature has Spearman $|\rho| > 0.8$ with any already-selected feature, it is discarded as redundant. This greedy de-duplication ensures the selected set captures distinct aspects of algorithm behaviour.

The final output is a set of 10 features, of which 3 come from the existing BLADE metrics and 7 are novel features introduced in this thesis, that are carried forward to Stage~3. The full selection results are presented in \cref{sec:results-features}.

% ---------------------------------------------------------------------------
\section{Feedback Format Screening}
\label{sec:meth-feedback-screening}

This stage tests whether adding a single behavioural feature to the LLM's feedback improves algorithm quality, and whether the \emph{framing} of that feedback matters.

\paragraph{Model.}
All Stage~3 experiments use Gemini~3 Flash~\cite{gemini3flash}, identified as the best model in Stage~1 on the composite of quality and efficiency (\cref{sec:results-screening}).

\paragraph{Experimental design.}
Each of the 10 selected features is tested individually. For each feature, three feedback formats are compared, yielding $10 \times 3 = 30$ conditions. One condition is excluded (comparative feedback for \texttt{longest\_no\_improvement\_streak}, due to a non-monotonic relationship with AOCC, as discussed below), giving 29 active conditions. Each condition is run for 100 candidates across 5 independent seeds (seeds 0--4), using the same shared initial population as Stage~1. The vanilla AOCC-only results from Stage~1's Gemini~3 Flash runs serve as the baseline.

\subsection{Feedback Formats}
\label{sec:meth-feedback-formats}

All three formats append a single behavioural feature to the standard AOCC feedback string. They differ in how much interpretive guidance they provide.

\paragraph{Neutral.}
Reports the feature value and its definition, without any indication of whether high or low values are desirable:

\begin{quote}
\small\itshape
The algorithm X got an average AOCC score of 0.7234 with standard deviation 0.0891. It achieved a \texttt{step\_size\_autocorrelation} of 0.8421 with standard deviation 0.0532, which measures whether large steps tend to follow large steps, capturing momentum in the search.
\end{quote}

\paragraph{Directional.}
Extends the neutral format with a sentence indicating which direction is associated with better performance. The direction is derived from the Stage~2 Spearman correlation analysis:

\begin{quote}
\small\itshape
\ldots capturing momentum in the search. Higher values are associated with better performance, indicating consistent, smooth step-size adaptation rather than erratic jumps.
\end{quote}

\paragraph{Comparative.}
Replaces the directional guidance with a distance-aware comparison against a concrete reference value. The reference is the median feature value of top-10\% algorithms by AOCC from the Stage~1 dataset ($n = 280$, AOCC $\geq 0.847$). The language adapts based on how far the candidate's value falls from this reference, using four tiers:

\begin{table}[t]
\centering
\caption{Comparative feedback tiers based on normalised distance from the top-10\% reference.}
\label{tab:comparative-tiers}
\footnotesize
\begin{tabular}{@{}rl@{}}
\toprule
\textbf{Distance} & \textbf{Tier} \\
\midrule
$\leq 0$      & Already meets reference \\
$0$--$0.15$   & Close to reference \\
$0.15$--$0.50$ & Moderate gap \\
$> 0.50$      & Far from reference \\
\bottomrule
\end{tabular}
\end{table}

The normalised distance is $d = (v_{\text{ref}} - v) / (v_{\text{top10}} - v_{\text{bot25}})$. For example, at distance $\leq 0$ the feedback reads: ``\textit{This already meets or exceeds the top-performing reference of 0.9070. Maintain this.}'' At distance $> 0.50$: ``\textit{This is far from the reference. Increasing this value is a priority.}''

The normalised distance is computed as $d = (v_{\text{ref}} - v) / (v_{\text{top10}} - v_{\text{bot25}})$, where $v$ is the candidate's feature value, $v_{\text{ref}} = v_{\text{top10}}$ is the top-10\% median, and $v_{\text{bot25}}$ is the bottom-25\% median. The bottom-25\% anchor (AOCC $\leq 0.233$, essentially RandomSearch-level performance) provides a meaningful scale: it represents the range between good and bad algorithms for each feature.

\paragraph{Excluded condition.}
Comparative feedback is excluded for \texttt{longest\_no\_improvement\_streak} because this feature has a U-shaped relationship with AOCC. Top-performing algorithms have high streak values ($\sim$5{,}327) because they converge early and coast through the remaining budget. Bottom-performing algorithms also have high values ($\sim$6{,}264) because they stagnate. A single reference target would be misleading in this non-monotonic case.

% ---------------------------------------------------------------------------
\section{Statistical Methodology}
\label{sec:meth-stats}

\paragraph{Overall format comparison.}
The Kruskal--Wallis test~\cite{kruskal1952ranks} is used to compare the distributions of best AOCC values across the three feedback formats (neutral, directional, comparative), pooling all features within each format.

\paragraph{Pairwise comparisons.}
Mann--Whitney $U$ tests~\cite{mann1947test} are applied for pairwise format comparisons within each feature, with Bonferroni correction for multiple comparisons.

\paragraph{Effect sizes.}
Cliff's delta~\cite{cliff1993dominance} is reported for all pairwise comparisons, with thresholds $|d| < 0.147$ (negligible), $< 0.33$ (small), $< 0.474$ (medium), $\geq 0.474$ (large).

\paragraph{AOCC-matched analysis.}
To disentangle whether feedback changes algorithm behaviour from whether it improves performance, we compare behavioural feature distributions \emph{within} AOCC quality bands. All valid candidates across all conditions are binned into four bands: low (AOCC 0.3--0.5), mid (0.5--0.7), high (0.7--0.85), and elite ($> 0.85$). Within each band, Kruskal--Wallis tests compare feature distributions across formats. This controls for algorithm quality and isolates the behavioural effect of feedback format.

% ---------------------------------------------------------------------------
\section{Infrastructure}
\label{sec:meth-infrastructure}

Experiments are executed on two LIACS compute servers:
\begin{itemize}
  \item \textbf{vibranium}: 2$\times$ NVIDIA RTX~3090 GPUs, used for serving larger Ollama models (27B, 32B).
  \item \textbf{duranium}: 6$\times$ NVIDIA GTX~980~Ti and 2$\times$ Titan~X GPUs, used for smaller models and parallel runs.
\end{itemize}
Local models are served via Ollama~\cite{ollama}, with vLLM~\cite{kwon2023vllm} available as an alternative backend for higher-throughput inference. Gemini~3 models~\cite{gemini, gemini3flash} are accessed through Google's Gemini API. All code, configuration, and analysis notebooks are available in the thesis repository.\footnote{\url{https://github.com/maxharell/thesis}} The BLADE benchmarking framework~\cite{vanStein2025BLADE} and LLaMEA~\cite{vanStein2024LLaMEA} are included as Git submodules, with minor modifications to support custom feedback strings and extended behavioural metric logging.
